\documentclass[11pt,a4paper]{article}
\usepackage[utf8]{inputenc}
\usepackage[T1]{fontenc}
\usepackage[portuguese]{babel}
\usepackage{xcolor}
\usepackage{hyperref}
\usepackage{geometry}
\usepackage{listings}
\usepackage{tcolorbox}
\geometry{margin=2.5cm}
\definecolor{primary}{RGB}{41,128,185}
\definecolor{secondary}{RGB}{44,62,80}
\definecolor{codebg}{RGB}{240,240,240}
\lstset{backgroundcolor=\color{codebg},basicstyle=\ttfamily\small,breaklines=true,frame=single}
\title{\color{secondary}\Huge\textbf{Organização de Cluster com Namespaces}}
\author{\color{primary}DevOps com Kubernetes -- 2026}
\date{}
\begin{document}
\maketitle


\section*{\color{primary}O que você aprenderá nesta página}
\begin{itemize}
\item Organizar recursos por namespace
\item Separar ambientes e responsabilidades
\item Aplicar comandos com contexto correto
\end{itemize}

Conforme a quantidade de aplicações cresce, colocar tudo no namespace \texttt{default} torna o cluster confuso. Namespaces são uma forma de separar recursos logicamente dentro do mesmo cluster. Eles não criam isolamento absoluto de segurança por si só, mas ajudam a organizar ambientes, equipes, aplicações e ferramentas de infraestrutura.

Um projeto pode usar namespaces como \texttt{dev}, \texttt{staging} e \texttt{prod}, ou separar por domínio, como \texttt{observability}, \texttt{database} e \texttt{apps}. O mais importante é que a separação reflita uma necessidade real de operação.

\begin{lstlisting}[language=bash]
kubectl create namespace dev
kubectl get namespaces
kubectl get pods -n dev
\end{lstlisting}

\section*{\color{primary}Aplicando recursos em um namespace}

Um recurso pode declarar o namespace no próprio manifesto ou receber o namespace no comando \texttt{kubectl}. Quando o namespace é declarado no YAML, o manifesto fica mais explícito. Quando ele é passado por linha de comando, o mesmo arquivo pode ser reutilizado em ambientes diferentes.

\begin{lstlisting}[language=bash]
apiVersion: v1
kind: Namespace
metadata:
  name: apps
---
apiVersion: apps/v1
kind: Deployment
metadata:
  name: api-dep
  namespace: apps
spec:
  replicas: 2
  selector:
    matchLabels:
      app: api
  template:
    metadata:
      labels:
        app: api
    spec:
      containers:
        - name: api
          image: nginx:1.27
\end{lstlisting}

\section*{\color{primary}Contextos do kubectl}

Também é possível configurar o namespace padrão de um contexto. Isso reduz repetição, mas exige cuidado: muitos erros em clusters reais acontecem porque alguém aplicou manifestos no namespace errado.

\begin{lstlisting}[language=bash]
kubectl config set-context --current --namespace=apps
kubectl config view --minify
\end{lstlisting}

\section*{\color{primary}Relação com configuração e acesso}

Namespaces delimitam onde vivem \texttt{ConfigMaps}, \texttt{Secrets}, \texttt{Services}, \texttt{Deployments}, \texttt{StatefulSets} e ferramentas como Prometheus. Um Service chamado \texttt{api-svc} no namespace \texttt{apps} não é o mesmo que um Service de mesmo nome no namespace \texttt{dev}. De outro namespace, a aplicação deve chamar \texttt{api-svc.apps} ou o nome completo do DNS interno.

\section*{\color{primary}Boas práticas}

Crie namespaces para separar responsabilidades reais. Use nomes simples, previsíveis e documentados. Em produção, combine namespaces com RBAC, quotas e políticas de rede.

\end{document}
